%\clearpage

\section {Introduction}
\label{sec:intro}

In this technical note, we report the results of a 3-flavour $\nu/\bar{\nu}$ joint analysis,
performed by the {\em VALOR} group,        
on the combined Run 1-6 dataset, using data taken in both 
neutrino mode (\TotalPotNu{} POT),
with the polarity set to forward horn current (FHC),
and antineutrino (\TotalPotNuBar{} POT) mode
with polarity set reverse horn current (RHC).
%This dataset was collected in two run periods:
%May - June 2014 (Run 5, \RunVPot{} POT in antineutrino beam mode) and November 2014 - June 2015 
%(Run 6, 
%\RunVIPot{} 
%POT in antineutrino beam mode).
This analysis builds on our other analyses presented in Refs. 
\cite{OurRun1DisappAnalysis},
\cite{OurRun12DisappAnalysis}, 
\cite{OurRun12DisappAnalysis3Flavour},
\cite{OurRun123DisappAnalysis3Flavour},
\cite{OurRun1234aDisappAnalysis3Flavour},
\cite{OurRun1234DisappAnalysis3Flavour},
\cite{OurRun123JointOscillationAnalysis},
\cite{OurRun1234JointOscillationAnalysis} 
\cite{OurRun56AntiNumuDisappearanceAnalysis3Flavour}
and
\cite{OurRun56AntiNueBarAppearanceAnalysis3Flavour}
As in \cite{OurRun56AntiNueBarAppearanceAnalysis3Flavour},
a hybrid frequentist-Bayesian method \cite{CousinsHighland}, 
which consists of marginalizing instead of profiling 
the nuisance parameters, is used.

The consistency of the data with the expectation is given by comparing the predicted reconstructed neutrino energy 
and outgoing lepton angle spectra of 1 ring events at Super-K to the observed spectra. 
The method of calculation of these predicted spectra is presented in Sec.~\ref{sec:spectrum_prediction}.
This analysis uses flux tuning based on recent NA61/SHINE 2009 data \cite{FluxTuning}\cite{NA61/SHINET2K}. 
It also uses an indirect extrapolation from ND280 to Super-K, which is implemented via a fit of flux and cross-section parameters to
ND280 measurements of the muon momentum-angle distributions of $\bar{\nu}_\mu$ and  $\nu_\mu$
%CC0$\pi$-like, CC1$\pi$-like and CC other events whose 
for different event topologies
with the neutrino interaction vertex either in FGD1 or FGD2 \cite{BANFF2016}.


% Ereco Vs Theta mu-like and e-like
A significant difference between this and the previous analyses is that
for $e$-like samples the analysis templates are built as a function of both 
the reconstructed neutrino energy and lepton angle,
in order to improve the separation of the different reaction / flavour components,
in particular $\bar{\nu}_{\mu} \rightarrow \bar{\nu}_e$ and $\nu_{\mu} \rightarrow \nu_e$,
fundamental in measuring $\delta_{CP}$.
Furthermore the $\nu_e$ intrinsic beam and NC components can also be better separated from signal.
Though the new variable $\theta$ increase the number of template bins and so the CPU time,
the number of reconstructed neutrino energy bins has not been reduced,
in order to continue to have a very good sensitivity to the oscillation pattern.
For $\mu$-like samples the templates are built as a function of the reconstructed neutrino energy
as in the past analyses.




Updated nominal predictions of the reconstructed neutrino energy and lepton angle spectrum 
of single ring events in Super-K are presented 
for Run 1-6 in Sec.~\ref{sec:spectrum_prediction}

The oscillation analysis is performed using neutrino oscillation probabilities computed in a 3-flavour
framework including constant-density matter effects. 
Details of our oscillation probability calculations, estimates of the numerical accuracy of our
calculations and comparisons with calculations used by the alternative T2K oscillation analyses 
are given in Appendices A and B of Ref. \cite{OurRun123DisappAnalysis3Flavour}.

As in our previous analyses, the oscillation fit is made using the likelihood-ratio method \cite{StatAnalysisCowan}.
The mathematical details of this method are given in Sec.\ 3.2 in Ref.\ \cite{OurRun1DisappAnalysis}.
We consider a total of \NSyst{} parameters including energy scale, neutrino flux, neutrino cross-section,
efficiency, initial-state nuclear environment and final-state re-scattering systematics 
plus the oscillation parameters (nuisance parameters).
In this Run 1-6 analysis, all the nuisance parameters are marginalised 
~ \cite{CousinsHighland}
.
The nuisance parameters and their effects 
on the extrapolated single ring event spectra are described in Sec.~\ref{sec:systematics}.
The analysis method is presented in detail in Sec.~\ref{sec:oscfit_method}

Fit validation and coverage studies are presented in Sec.~\ref{sec:validation}.

The main goal of this analysis is to search for CP violation in the context of the PMNS matrix
and confidence intervals are set using the Feldman-Cousins method \cite{FeldmanCousins}.
Confidence intervals of the other oscillation parameters are set using the 
constant $\Delta \chi^2$ method.
Sensitivity studies are shown in Sec.~\ref{sec:sensi}.
The results of the measurement using the Run 1-6 dataset is given in Sec.~\ref{sec:results} of this note. 

%In Sec.~\ref{sec:nuebarapp},~\ref{sec:sensitivity_nuebarapp} and~\ref{sec:results_nuebar} the $\bar{\nu}_e$ appearance analysis is shown 

%while the $\bar{\nu}_{\mu}$ disappearance analysis is shown in Sec.~\ref{sec:numubardis} and~\ref{sec:oscfit_numubardis}.
Finally, a summary is given in Sec.~\ref{sec:summary}.  

\clearpage
